\begin{titlepage}
  \setlength{\parindent}{0cm}

  \includegraphics[width=200px]{Images/uob-logo-black-transparent}

  \Large
  Faculty of Engineering and Design

  \vspace{80pt}

  \LARGE
  Final Year MEng Project \\
  Project Report

  \vspace{80pt}
  \textbf{Robot Football - High Level Control and Strategy Implementation}

  \vspace{10pt}
  \emph{Anthony Richards} \\
  \emph{\today}

  \vspace{80pt}
  Supervisor: \emph{Dr.~Pejman Iravani}

  \vspace{10pt}
  Assessor: \emph{TBD}
\end{titlepage}

\pagenumbering{roman}

\begin{abstract}
TODO : Abstract
\end{abstract}

\tableofcontents

\section*{Document Conventions}
The following unit conventions are used by this document

\begin{itemize}
  \item Unless otherwise stated, all units used are standard SI units
  \item The symbol 'in' is used to designate the unit inches
\end{itemize}

All images displaying field layouts are shown as a grayscale image representing
relative magnitude.  A white pixel represents a point of highest magnitude
(either gradient or field intensity), while a black pixel represents a point of
lowest magnitude.

\subsection*{List of Acronyms}

The following acronyms are in use in this document:

\begin{itemize}
  \item CPU - Central Processing Unit
  \item CRT - C Runtime Library
  \item CSV - Comma Seperated Variable
  \item DLL - Dynamic Link Library
  \item FIRA - Federation of International Robot-Soccor Association
  \item GPR - General Purpose Register
  \item GPU - Graphics Processing Unit
  \item SIMD - Single Instruction Multiple Data
\end{itemize}

\subsection*{List of Symbols}

The following symboles are in use in this document:

\begin{itemize}
  \item $e$ - The base of the natural logarithm.
  \item $\boldsymbol{F}(x,y)$ - The gradient of a potential field at point
  $\left(x,y\right)$.  The real part is the gradient in $x$, while the imaginary
  part is the gradient in $y$.
  \item $G(s)$ - A transfer function, represented in the laplace domain.
  \item $k_{attr}$ - A weight for an attractive field.
  \item $k_{repulse}$ - A weight for a repulsive field.
  \item $N$ - The number of elements in a summation.
  \item $p_i(x,y)$ - The intensity of potential field $i$ at point
  $\left(x,y\right)$
  \item $P(x,y)$ - The total potential field intensitity at point
  $\left(x,y\right)$.
  \item $R_0$ - The limiting radius of a field.
  \item $R^2$ - A measure of the quality of a line of best fit, where 1 is a
  perfect fit a smaller number represents a worse fit.
  \item $t$ - The time in seconds (normally relative to the beginning of the
  operation).
  \item $\sigma$ - The variable controlling the spread of a Gaussian function
\end{itemize}

\cleardoublepage

\pagenumbering{arabic}